% SPDX-License-Identifier: CC-BY-4.0
\documentclass[11pt,letterpaper]{article}

\usepackage[T1]{fontenc}
\usepackage{amsmath,amssymb}
\usepackage{tikz}
\usepackage{pgfplots}
\pgfplotsset{compat=1.18}
\usepackage{algorithm}
\usepackage{algpseudocode}
\usepackage{booktabs}
\usepackage{enumitem}
\usepackage[hidelinks]{hyperref}

\newcommand{\N}{\mathbb{N}}
\newcommand{\Z}{\mathbb{Z}}
\newcommand{\Q}{\mathbb{Q}}
\newcommand{\R}{\mathbb{R}}
\newcommand{\C}{\mathbb{C}}
\newcommand{\Map}[3]{#1 \colon #2 \to #3}
\newcommand{\Mod}[1]{#1\text{-}\mathrm{Mod}}
\newcommand{\GroupHom}{\mathrm{Hom}_{\mathbf{Grp}}}
\newcommand{\RingHom}{\mathrm{Hom}_{\mathbf{Ring}}}

\newtheorem{theorem}{Theorem}
\newtheorem{lemma}[theorem]{Lemma}
\newtheorem{corollary}[theorem]{Corollary}
\newtheorem{definition}[theorem]{Definition}
\newtheorem{remark}[theorem]{Remark}
\newenvironment{proof}{\noindent\textit{Proof.}}{\hfill\(\square\)}

\title{Generalized Collatz Library: Exact Statistics of the Raw Dyadic Clock}
\author{}
\date{\today}

\begin{document}
\maketitle

\begin{abstract}
For odd multiplier and offset on the 2-adic integers, raw affine events
cannot occur consecutively. Their parity process is a two-state Markov
chain. We derive exact finite word counts and moments, a central limit
theorem, and an explicit large-deviation rate. Cylinder probabilities
transfer to ordinary integer densities with a uniform discrepancy
bound. The same bound gives a central limit for integer samples at
growing finite depths; almost-sure statements concern the 2-adic
probability space.
\end{abstract}

\section{The probability space and the raw clock}

Let \(\Z_2\) denote the ring of 2-adic integers. Fix odd
\(q,c\in\Z_2\), and use
\[
T(x)=\begin{cases}x/2,&x\text{ even},\\qx+c,&x\text{ odd}.\end{cases}
\]
Write \(\mu\) for normalized Haar probability and \(\nu\) for the
invariant probability with Haar density \(4/3\) on even states and
\(2/3\) on odd states. For \(h\ge1\), let
\[
S_j(x)=\boldsymbol{1}_{\{T^j(x)\text{ odd}\}},\qquad
B_h(x)=\sum_{j=0}^{h-1}S_j(x).
\]
Thus \(B_h\) counts affine instructions in \(h\) raw ticks, and
\(h-B_h\) counts divisions. All logarithms are natural.

The companion cylinder construction \cite{DyadicTopology26} proves
that the parity sequence has no adjacent ones. If a legal length-\(h\)
word has \(b\) ones and last letter \(\delta\in\{0,1\}\), its Haar
measure is \(2^{-h+b-\delta}\). Under \(\nu\) its measure is this
quantity multiplied by \(4/3\) or \(2/3\), according as its first
letter is zero or one. The parity process is a Markov chain with
\[
P=\begin{pmatrix}1/2&1/2\\1&0\end{pmatrix}.
\]
Its initial distribution is \((1/2,1/2)\) under \(\mu\), and the
stationary distribution \((2/3,1/3)\) under \(\nu\).
Indeed, conditioned on an even source, the next state is Haar
distributed; conditioned on an odd source, it is Haar distributed
within the even states. Applying this argument successively proves
the conditional transition law, not just its one-step marginal.

\section{Exact counts and moments}

% Removing one separating zero after each earlier one gives a finite bijection.
\begin{lemma}[Word counts]
For \(0\le b\le\lceil h/2\rceil\), the number of legal length-\(h\)
words with \(b\) ones is
\[
N(h,b)=\binom{h-b+1}{b}.
\]
Consequently,
\[
\begin{aligned}
\tfrac12 N(h,b)2^{-h+b}
&\le\mu(B_h=b)\le N(h,b)2^{-h+b},\\
\tfrac13 N(h,b)2^{-h+b}
&\le\nu(B_h=b)\le\tfrac43 N(h,b)2^{-h+b}.
\end{aligned}
\]
Outside that range both probabilities are zero.
\end{lemma}

\begin{proof}
If the one positions are \(1\le i_1<\cdots<i_b\le h\), legality says
\(i_{r+1}\ge i_r+2\). The positions \(j_r=i_r-(r-1)\) form an
arbitrary \(b\)-element subset of \(\{1,\ldots,h-b+1\}\).
The inverse is \(i_r=j_r+(r-1)\), so this is a bijection, also for
\(b=0\). Each Haar word weight differs from \(2^{-h+b}\) by the
factor \(2^{-\delta}\in\{1,1/2\}\). Multiplying by the specified
first-letter density gives factors between \(1/3\) and \(4/3\)
for \(\nu\). Summing over the bijection proves the bounds.
\end{proof}

% The nontrivial Markov eigenvalue controls every covariance and the finite variance.
\begin{theorem}[Exact stationary moments]
\label{thm:moments}
For every \(h\ge1\),
\[
\mathbb E_\nu B_h=h/3,\qquad
\operatorname{Var}_\nu B_h
=\frac{2h}{27}+\frac8{81}\bigl(1-(-1/2)^h\bigr).
\]
For Haar initial data,
\[
\mathbb E_\mu B_h=h/3+\frac19\bigl(1-(-1/2)^h\bigr).
\]
\end{theorem}

\begin{proof}
Let \(\Pi\) have both rows equal to \((2/3,1/3)\). Direct
multiplication gives \(P^r=\Pi+(-1/2)^r(I-\Pi)\).
Under the stationary law,
\[
\operatorname{Cov}_\nu(S_0,S_r)=\frac29(-1/2)^r.
\]
This includes \(r=0\), since a Bernoulli variable with mean \(1/3\)
has variance \(2/9\). Summing the covariances gives
\[
\operatorname{Var}_\nu B_h
=\frac{2h}9+\frac49\sum_{r=1}^{h-1}(h-r)(-1/2)^r.
\]
For \(t\ne1\), multiplication of the finite geometric sum by
\(1-t\), or differentiation of that polynomial, yields
\[
\sum_{r=1}^{h-1}(h-r)t^r
=\frac{ht}{1-t}-\frac{t}{(1-t)^2}
+\frac{t^{h+1}}{(1-t)^2}.
\]
Substitution of \(t=-1/2\) proves the variance formula.
Under \(\mu\), the odd probability at time \(r\) is
\(1/3+(1/6)(-1/2)^r\), by the same matrix identity and the initial
law \((1/2,1/2)\). Its sum proves the final assertion.
\end{proof}

\section{Large deviations}

Define, with \(0\log0=0\),
\[
I(a)=
\begin{cases}
\begin{aligned}
&(1-a)\log2+a\log a\\
&\quad +(1-2a)\log(1-2a)-(1-a)\log(1-a)
\end{aligned}
&0\le a\le1/2,\\
+\infty,&\text{otherwise}.
\end{cases}
\]
It is continuous on its finite domain, with \(I(0)=\log2\),
\(I(1/2)=\tfrac12\log2\), and its unique zero is \(a=1/3\).
To check the last assertion directly, on \(0<a<1/2\),
\[
I'(a)=\log\frac{a(1-a)}{2(1-2a)^2},\qquad
I''(a)=\frac1a-\frac1{1-a}+\frac4{1-2a}>0.
\]
The derivative vanishes at \(1/3\), where substitution gives zero.

% Uniform factorial estimates turn the exact finite count into both LDP directions.
\begin{theorem}[Explicit rate]
\label{thm:rate}
For either \(\eta=\mu\) or \(\eta=\nu\), the variables \(B_h/h\)
satisfy the large-deviation principle with rate \(I\). Precisely,
for every closed \(F\subset\R\) and open \(U\subset\R\),
\[
\begin{aligned}
\limsup_{h\to\infty}h^{-1}\log\eta(B_h/h\in F)
&\le-\inf_{a\in F}I(a),\\
\liminf_{h\to\infty}h^{-1}\log\eta(B_h/h\in U)
&\ge-\inf_{a\in U}I(a).
\end{aligned}
\]
The convention is \(\log0=-\infty\).
\end{theorem}

\begin{proof}
Put \(f(m)=m\log m\) for positive integers and \(f(0)=0\).
Comparison with the integral of the increasing function \(\log x\)
gives \(\log(m!)=f(m)-m+O(\log(m+1))\), with an absolute constant.
Therefore, uniformly for legal \(b\),
\[
\log N(h,b)
=f(h-b+1)-f(b)-f(h-2b+1)+O(\log(h+1)).
\]
Replacing an argument \(m+1\) by \(m\) changes \(f\) by at most
\(\log(h+1)+1\). For \(b\le h/2\), apply these two replacements,
then add \((-h+b)\log2\) and the bounded word-weight factors.
This yields
\[
\log\eta(B_h=b)=-hI(b/h)+O(\log(h+1)).
\]
If \(h\) is odd and \(b=(h+1)/2\), then \(N(h,b)=1\);
the same estimate holds with \(b/h\) replaced by \(1/2\).
Thus the uniform formula uses
\(\widehat a_h(b)=\min(b/h,1/2)\), with
\(\lvert\widehat a_h(b)-b/h\rvert\le1/(2h)\).

For a closed set, sum at most \(h+1\) probabilities and use this
uniform formula. Any convergent sequence of admissible \(b/h\) in
the closed set has its limit in that set and in \([0,1/2]\).
Compactness and continuity of \(I\) prove the upper bound.
If there is no such sequence, the event is eventually empty.

For an open set and any \(a\in U\cap[0,1/2]\), choose legal
integers \(b_h\) with \(b_h/h\to a\). These exist at both endpoints:
take zero for \(a=0\), and \(\lfloor h/2\rfloor\) for \(a=1/2\).
For all large \(h\), \(b_h/h\in U\). A single point probability
and the uniform formula give the lower bound \(-I(a)\).
Take the supremum of these lower bounds; an empty intersection
makes the required lower bound vacuous.
\end{proof}

\section{The central limit and almost-sure frequency}

% The analytic two-eigenvalue formula supplies a characteristic-function limit.
\begin{theorem}[Raw central limit]
\label{thm:clt}
Under either \(\mu\) or \(\nu\),
\[
\frac{B_h-h/3}{\sqrt{2h/27}}
\ \Longrightarrow\ N(0,1).
\]
\end{theorem}

\begin{proof}
For either initial row \(p=(1/2,1/2)\) or \(p=(2/3,1/3)\), put
\(D(s)=\operatorname{diag}(1,e^s)\). Conditioning successively on
the parity states gives the exact moment generating function
\[
\mathbb E_p e^{sB_h}=pD(s)\bigl(PD(s)\bigr)^{h-1}
\begin{pmatrix}1\\1\end{pmatrix}.
\]
The eigenvalues of \(PD(s)\), analytic near zero with the indicated
square-root branch, are
\[
\lambda(s)=\frac{1+\sqrt{1+8e^s}}4,\qquad
\gamma(s)=\frac{1-\sqrt{1+8e^s}}4.
\]
They are distinct near zero. The two spectral projections give
\(\mathbb E_p e^{sB_h}
=A_p(s)\lambda(s)^{h-1}+C_p(s)\gamma(s)^{h-1}\),
where the coefficients are analytic. At zero, the eigenvector
\((1,1)^\mathsf{T}\) lies entirely in the eigenvalue-one space;
thus \(A_p(0)=1\), \(C_p(0)=0\).
Moreover \(\lambda(0)=1\), \(\gamma(0)=-1/2\), so the second
term is exponentially small relative to the first, uniformly in
a sufficiently small complex neighborhood of zero.

Differentiation of the displayed root gives
\[
\log\lambda(s)=s/3+s^2/27+O(|s|^3).
\]
For fixed real \(t\), set \(s=it/\sqrt{2h/27}\), and subtract the
centering term \(sh/3\). The logarithm of the leading term tends
to \(-t^2/2\); its cubic remainder is \(O(h^{-1/2})\).
The analytic prefactor tends to one, and the other eigenvalue
contribution tends to zero. Hence the characteristic functions
of the normalized variables converge to \(e^{-t^2/2}\).

Here is an explicit passage from this limit to distribution functions.
The stationary variance formula bounds the second moments of the
normalized variables uniformly. This is also true under \(\mu\),
because \(d\mu/d\nu\le3/2\). Add an independent centered Gaussian
\(Z_\varepsilon\) of variance \(\varepsilon^2>0\).
Gaussian Fourier inversion and Fubini give the smoothed density
\[
g_{h,\varepsilon}(u)=\frac1{2\pi}\int_\R
e^{-itu}\phi_h(t)e^{-\varepsilon^2t^2/2}\,dt.
\]
Dominated convergence implies uniform convergence of these
densities to the centered Gaussian density of variance
\(1+\varepsilon^2\). Uniform second-moment bounds control the
tails, so integration on bounded intervals proves convergence
of the smoothed distribution functions.
For every \(\delta>0\), the unsmoothed distribution function
\(F_h\) lies between
\[
F_{h,\varepsilon}(u-\delta)-\Pr(|Z_\varepsilon|>\delta)
\quad\text{and}\quad
F_{h,\varepsilon}(u+\delta)+\Pr(|Z_\varepsilon|>\delta).
\]
First let \(h\to\infty\). Then take \(\delta=\sqrt\varepsilon\)
and let \(\varepsilon\downarrow0\); Chebyshev bounds the extra
tail by \(\varepsilon\), and the Gaussian distribution function
is continuous. These bounds prove the asserted central limit.
\end{proof}

% A summable subsequence estimate and bounded increments recover the full sequence.
\begin{corollary}[2-adic frequency]
For \(\mu\)-almost every and \(\nu\)-almost every \(x\),
\[
\lim_{h\to\infty}B_h(x)/h=1/3.
\]
The division frequency is therefore \(2/3\). The 2-adic derivative
Lyapunov exponent is \((2/3)\log2\).
\end{corollary}

\begin{proof}
The stationary variance is \(O(h)\). At \(h=k^2\), Chebyshev gives
\(\nu(|B_{k^2}/k^2-1/3|>\varepsilon)=O_\varepsilon(k^{-2})\).
The probability of any such event after index \(K\) is bounded
by the tail of this convergent series, which tends to zero.
Using \(\varepsilon=1/m\), \(m\ge1\), gives almost-sure convergence
along squares. For \(k^2\le h<(k+1)^2\), the count changes by at
most \(2k+1\), proving full-sequence convergence. The measures
\(\mu,\nu\) are equivalent, so the same almost-sure conclusion
holds for Haar data.

On each parity neighborhood the derivative is \(1/2\) or \(q\).
The chain rule is applicable on each finite itinerary cylinder.
Since \(|q|_2=1\), it gives
\(\lvert(T^h)'(x)\rvert_2=2^{h-B_h(x)}\).
Taking its logarithm and using the proved frequency gives
the exponent.
\end{proof}

For ordinary nonzero odd integer \(q\), the associated
\emph{real slope statistic} has limiting value
\((\log|q|-2\log2)/3\) on this probability space. This statistic
is the logarithm of the absolute product of the branch slopes,
divided by the raw time. It is not a claim about the growth of ordinary
integer orbit values.

\section{Recovery of finite integer laws}

% Finite precision gives a uniform discrepancy, while the stated limit order excludes capture.
\begin{theorem}[Ordinary density transfer]
\label{thm:transfer}
Suppose now that \(q,c\) are odd integers. For any event \(E\)
depending only on the first \(h\ge1\) raw parity letters,
\[
\left|\frac{\#\{1\le n\le X:n\in E\}}X-\mu(E)\right|
\le\frac{2^h}{X}
\]
for every positive integer \(X\).
Thus every fixed finite-count event has ordinary natural density
given by its Haar law. The central limit also holds for uniformly
selected \(n\in\{1,\ldots,X\}\) at any integer depth \(h(X)\to\infty\)
with \(2^{h(X)}/X\to0\).
\end{theorem}

\begin{proof}
Each legal word cylinder is one residue class modulo
\(2^{h-b+\delta}\), and \(h-b+\delta\le h\). Thus \(E\) is a union
of classes modulo \(2^h\). Each class contributes within one of
\(X/2^h\), and its Haar measure is \(2^{-h}\).
Summing these discrepancies over at most \(2^h\) classes proves
the bound. For a distribution-function event in the central limit,
this error is uniform in the threshold. Combine it with the Haar
central limit at the growing depth to prove the last assertion.
\end{proof}

These recovery statements concern finite prefixes. For the
large-deviation law, one valid interpretation is to take the
ordinary-density limit \(X\to\infty\) at each fixed \(h\), then
let \(h\to\infty\). The uniform bound also controls other joint
limits when its error is small compared with the probability
being estimated. None of these conclusions alone assigns an
infinite-orbit frequency or capture outcome to density-one
ordinary integer seeds.

% Integer capture and 2-adic typical frequencies can differ even for the same coefficients.
\begin{remark}[An integer comparison]
For \(q=-1,c=1\), every ordinary integer reaches zero, so
\(B_h(n)/h\to0\) for every integer seed \(n\). To prove this,
use the first-division map
\[
H(n)=\begin{cases}n/2,&n\text{ even},\\(1-n)/2,&n\text{ odd}.
\end{cases}
\]
For \(|n|>1\), its magnitude is strictly smaller than \(|n|\):
the even case halves it, and the odd case satisfies
\(|H(n)|\le(1+|n|)/2<|n|\). Descent in the nonnegative
integer magnitude reaches \(\{-1,0,1\}\); here
\(H(-1)=1\), \(H(1)=0\), and \(H(0)=0\).
Every \(H\)-step is one or two legal raw ticks, so the raw orbit
also reaches zero. Its total number of odd visits is then finite.
Nevertheless, the almost-sure 2-adic limit for these same coefficients
is \(1/3\), by the frequency corollary. This example separates
fixed-prefix integer density from infinite integer-orbit frequency.
\end{remark}

\section{Exact finite implementation}

The accompanying Go implementation \cite{DyadicCode26} enumerates every
legal word through depth sixteen, compares counts with the
binomial formula, and verifies stationary normalization, mean
and variance by exact integer arithmetic. It also checks
the raw cylinder formulas against signed orbits for odd integer
\(q,c\in[-7,7]\) through depth twelve. The first computation tests
Theorem~\ref{thm:moments}; the second tests the cylinder input to
Theorem~\ref{thm:transfer}. Neither finite enumeration supplies the
infinite arguments in Theorems~\ref{thm:rate} and~\ref{thm:clt},
which are proved above. The statistical theorems in this article
are written proofs rather than Lean declarations.

\bibliographystyle{amsalpha}
\bibliography{DyadicReferences}
\end{document}
